../cictikz/src/cictikz/data/tex/cictikz_preamble.tex

% Nauta's transconductor. Everything in it is an inverter, which is why
% it works at any supply where an inverter still has gain.
%
% Laid out in four columns so the reader can see which inverter does
% what: the two amplifying inverters on the left, then the cross
% coupled pair between the output rails - each one taking the other
% rail and driving this one - and finally a shorted inverter hanging
% off each output as its load. Drawing the small inverters in their
% own columns, rather than on the rails, is what keeps the picture
% readable.

\begin{circuitikz}[american, thick, transform shape, circuit ee IEC]

  ../cictikz/src/cictikz/data/tex/cictikz_lib.tex

  \def\yt{2.3}       % top output rail
  \def\yb{-2.3}      % bottom output rail
  \def\xout{9.6}     % where the rails end

  % ---- the two amplifying inverters ----
  \draw (-2.4,\yt) \cicInv{inv1};
  \draw (-3.9,\yt) node[anchor=east] {$v_p$} to [short,o-] (inv1_in);
  \draw (-2.4,\yb) \cicInv{inv2};
  \draw (-3.9,\yb) node[anchor=east] {$v_n$} to [short,o-] (inv2_in);

  % ---- the output rails ----
  \draw (inv1_out) -- (\xout,\yt);
  \draw (\xout,\yt) to [short,-o] ++(0.5,0) node[anchor=west] {$v_{on}$};
  \draw (inv2_out) -- (\xout,\yb);
  \draw (\xout,\yb) to [short,-o] ++(0.5,0) node[anchor=west] {$v_{op}$};

  % ---- the cross coupled pair, each in its own column ----
  % top rail drives the bottom one
  \begin{scope}[shift={(1.1,0.8)}, rotate=-90, scale=0.8]
    \draw (0,0) \cicInv{inv3};
  \end{scope}
  \draw (1.1,\yt) -- (1.1,0.8);
  \fill (1.1,\yt) circle (0.075);
  \draw (inv3_out) -- (1.1,\yb);
  \fill (1.1,\yb) circle (0.075);

  % bottom rail drives the top one
  \begin{scope}[shift={(3.1,-0.8)}, rotate=90, scale=0.8]
    \draw (0,0) \cicInv{inv4};
  \end{scope}
  \draw (3.1,\yb) -- (3.1,-0.8);
  \fill (3.1,\yb) circle (0.075);
  \draw (inv4_out) -- (3.1,\yt);
  \fill (3.1,\yt) circle (0.075);

  %- the annotation goes below the picture, not between the devices
  \draw[gray!60, thin] (2.1,{\yb-0.55}) -- (1.1,{\yb-0.15});
  \draw[gray!60, thin] (2.1,{\yb-0.55}) -- (3.1,{\yb-0.15});
  \node[anchor=north, scale=0.85, align=center] at (2.1,{\yb-0.6})
    {cross coupled: each rail drives the other};

  % ---- the loads: a shorted inverter hanging off each output ----
  \draw (6.3,\yt) -- (6.3,{\yt+1.15});
  \fill (6.3,\yt) circle (0.075);
  \begin{scope}[shift={(6.3,{\yt+1.15})}, scale=0.8]
    \draw (0,0) \cicInv{inv5};
  \end{scope}
  \draw (inv5_out) -- (8.3,{\yt+1.15}) -- (8.3,\yt);
  \fill (8.3,\yt) circle (0.075);

  \draw (6.3,\yb) -- (6.3,{\yb-1.15});
  \fill (6.3,\yb) circle (0.075);
  \begin{scope}[shift={(6.3,{\yb-1.15})}, scale=0.8]
    \draw (0,0) \cicInv{inv6};
  \end{scope}
  \draw (inv6_out) -- (8.3,{\yb-1.15}) -- (8.3,\yb);
  \fill (8.3,\yb) circle (0.075);

  \node[anchor=south, scale=0.85] at (7.3,{\yt+1.95})
    {shorted inverters: the load};

\end{circuitikz}
\end{document}
